% =============================================================================
% 4.2 Experiment 2 Results — Kubernetes Isolation (SQ2)
% =============================================================================

\section{Experiment 2 Results --- Kubernetes Isolation (SQ2)}
\label{sec:results-exp2}

\subsection{Part A --- Isolation Comparison}
\label{subsec:results-exp2a}

Experiment~2A ran the same two-phase \texttt{thesis\_workload} job on the Kind cluster under the K8sRunLauncher.
Each job executes as an isolated pod with limits of 1\,vCPU and 2\,GiB RAM (the 400\,MB workload buffer plus $\approx$500--600\,MB Python/Dagster overhead requires at least 1.8\,GiB peak RSS).
Table~\ref{tab:results-k8s} shows that Kubernetes maintained a 100\% success rate across all concurrency levels.

\begin{table}[H]
  \centering
  \caption{Kubernetes deployment results by workload level (mean $\pm$ std over 3 repetitions, Kind single-node cluster, 4\,vCPU / 8\,GB node, 2\,GiB per-pod memory limit).
  All concurrency levels maintained 100\% success rate under stable node-level conditions.}
  \label{tab:results-k8s}
  \small
  \begin{tabular}{clccc}
    \toprule
    \textbf{Level} & \textbf{Concurrent Jobs} & \textbf{Success (\%)} & \textbf{Mean time (s)} & \textbf{Std time (s)} \\
    \midrule
    L1  & 1  & 100.0 & 63.4 & 0.1 \\
    L2  & 2  & 100.0 & 64.1 & 0.6 \\
    L3  & 3  & 100.0 & 64.6 & 0.5 \\
    L4  & 5  & 100.0 & 68.2 & 2.3 \\
    L5  & 7  & 100.0 & 80.6 & 9.6 \\
    L6  & 10 & 100.0 & 81.6 & 8.0 \\
    \bottomrule
  \end{tabular}
\end{table}

\begin{figure}[H]
  \centering
  \includegraphics[width=0.85\textwidth]{k8s_dashboard.png}
  \caption{Kubernetes experiment dashboard: success rate, execution time distribution, and VM vs.\ K8s comparison (Exp2A, $n=3$ reps). Success rate remains 100\,\% at all levels; execution time grows modestly from 63\,s to 82\,s, reflecting effective pod-level resource isolation.}
  \label{fig:k8s-isolation}
\end{figure}

\begin{figure}[H]
  \centering
  \includegraphics[width=0.85\textwidth]{sq2_vm_vs_k8s_comparison.png}
  \caption{VM vs.\ Kubernetes comparison: success rate (left) and execution time (right) by concurrency level (Exp2A, $n=3$ reps). K8s maintains 100\,\% success at all levels while VM reliability collapses at L4+. Shaded bands show $\pm$1 standard deviation.}
  \label{fig:k8s-success-rate}
\end{figure}

\begin{figure}[H]
  \centering
  \includegraphics[width=0.85\textwidth]{k8s_exec_time_distribution.png}
  \caption{Kubernetes job execution time distribution by concurrency level (Exp2A, $n=3$ reps). Boxes show interquartile range; whiskers extend to 1.5$\times$~IQR. The distribution widens at L5--L6 as node-level memory pressure increases, but remains tightly bounded at lower levels.}
  \label{fig:k8s-exec-dist}
\end{figure}

\paragraph{Key observation.}
Kubernetes maintained 100\,\% success rate across all tested levels (L1--L6) with stable execution times.
Mean execution time increased modestly from 63.4\,s (L1) to 81.6\,s (L6), a relative increase of $\approx$29\,\%.
The standard deviation at L5 (9.6\,s) and L6 (8.0\,s) is higher than at lower levels, reflecting mild aggregate memory pressure at 7--10 concurrent pods on the 8\,GB node, but remaining well within the pod's 2\,GiB limit.
Execution time variance remained low and stable across all levels, confirming that pod-level resource limits effectively shielded jobs from inter-job contention throughout the tested concurrency range.

\begin{figure}[H]
  \centering
  \includegraphics[width=0.85\textwidth]{sq4_crossover_point.png}
  \caption{Crossover analysis: cost per successful job delivery (left) and reliability divergence (right) vs.\ concurrency level. Left: VM effective time accounts for failed runs; K8s total includes startup overhead. Right: VM success rate collapses at L4 while K8s maintains 100\,\%. Reliability crossover at L4 (5 concurrent jobs).}
  \label{fig:vm-vs-k8s}
\end{figure}

\subsection{Part B --- Blast Radius Test}
\label{subsec:results-exp2b}

The blast radius test was conducted at level L5 (5 concurrent jobs).
After all five jobs were confirmed running, one job was deliberately terminated.
On Kubernetes, the target pod was deleted using \texttt{kubectl delete pod --force}.
On the VM, the equivalent Docker container process was sent \texttt{SIGKILL}.

\begin{table}[H]
  \centering
  \caption{Blast radius comparison at L5 (5 concurrent jobs) --- deliberate single-job failure injection.}
  \label{tab:blast-radius}
  \small
  \begin{tabular}{lcc}
    \toprule
    \textbf{Metric} & \textbf{VM (\texttt{kill -9})} & \textbf{K8s (\texttt{pod delete})} \\
    \midrule
    Failed job status & FAILURE (SIGKILL) & FAILURE (pod deleted) \\
    Other jobs affected (\#) & 0 & 0 \\
    Surviving jobs success rate (\%) & 100 & 100 \\
    Surviving jobs mean time impact (s) & $<$2 & $<$1 \\
    \bottomrule
  \end{tabular}
\end{table}

Both environments demonstrated full blast radius containment at this concurrency level:
the surviving four jobs completed successfully regardless of whether the failure occurred on the VM or in Kubernetes.
This result is consistent with the Docker isolation model on the VM where each job runs in its own container namespace, and with the Kubernetes pod isolation model.
The difference between environments becomes apparent at higher concurrency levels where VM memory exhaustion can cause system-level OOM that terminates multiple containers simultaneously (Section~\ref{subsec:exp1-degradation}).

\subsection{Part C --- Scheduling Overhead (SQ3 Supplementary)}
\label{subsec:results-exp2c}

Pod scheduling latency and container startup time were collected from Kubernetes pod event timestamps at each concurrency level (Table~\ref{tab:overhead-summary}).
These data directly address SQ3 and feed into the crossover analysis in Section~\ref{sec:results-exp3}.

\paragraph{Measurement methodology.}
Startup overhead is computed by joining \texttt{pod\_timing.csv} with \texttt{dagster\_runs.csv} on \texttt{run\_id}: startup\,overhead = total\,pod\,lifetime $-$ Dagster\,exec\,time.
Note: the \texttt{job\_start\_ts} field records the \texttt{ContainersReady} condition change on container \emph{exit}, so subtracting \texttt{submitted\_ts} gives total pod lifetime, not startup time.

At L1--L6, all pods were scheduled within $<$1\,s of submission and container startup overhead measured 24.4\,s (L1) to 56.6\,s (L6), growing with concurrency as more containers compete for Python interpreter init and PostgreSQL handshake slots.
No scheduling timeouts or pending states were observed at any tested level (Table~\ref{tab:overhead-summary}).
